\documentclass[aspectratio=169]{beamer}
\usetheme{Madrid}
\usecolortheme{default}
\usepackage[utf8]{inputenc}
\usepackage[T5]{fontenc}
\usepackage{graphicx}
\usepackage{hyperref}
\usepackage{booktabs}

\setbeamertemplate{caption}[numbered]
\renewcommand{\figurename}{Hình}

\title[Buổi 11]{Trí tuệ Nhân tạo cho Kế toán \\ \vspace{0.3cm} \Large Ứng dụng AI trong Kế toán Thuế}
\author{Đại học Đông Á}
\date{\today}

\begin{document}

% SLIDE 1
\begin{frame}
    \titlepage
\end{frame}

% SLIDE 2
\begin{frame}{Nội dung Bài học}
    \tableofcontents
\end{frame}


\section{BUỔI 11: ỨNG DỤNG AI TRONG KẾ TOÁN THUẾ \& TUÂN THỦ PHÁP LUẬT}


\begin{frame}{Năng lực đạt được sau buổi học}
    \begin{itemize}
        \item \textbf{Hiểu và ứng dụng AI:} Trong các quy trình kê khai thuế (GTGT, TNCN) và tự động hóa tính toán thuế.
        \item \textbf{Phân tích văn bản pháp luật:} Vận dụng AI để đọc hiểu, tóm tắt các thông tư, nghị định phức tạp về thuế (ví dụ: Thông tư 200, Luật Quản lý Thuế).
        \item \textbf{Quản lý rủi ro:} Sử dụng AI để nhận diện, đánh giá và kiểm soát các rủi ro pháp lý và rủi ro thuế từ các giao dịch.
        \item \textbf{Trách nhiệm nghề nghiệp:} Nắm bắt được giới hạn pháp lý và trách nhiệm đạo đức nghề nghiệp khi sử dụng AI.
    \end{itemize}
\end{frame}


\begin{frame}{Nội dung chương trình}
    1. Tổng quan Kế toán Thuế trong Kỷ nguyên số\\
    2. Tối ưu hóa Lập kế hoạch Thuế với AI (AI-Based Tax Planning)\\
    3. Ứng dụng AI Hỗ trợ Kê khai \& Tuân thủ Thuế\\
    4. Phát hiện Gian lận \& Quản trị Rủi ro Thuế\\
    5. Ứng dụng AI với Luật Quản lý Thuế \& Thông tư 200\\
    6. Giới hạn Pháp lý \& Đạo đức Nghề nghiệp\\
    7. Tổng kết \& Thảo luận\\
\end{frame}


\section{PHẦN 1: TỔNG QUAN KẾ TOÁN THUẾ TRONG KỶ NGUYÊN SỐ}


\begin{frame}{1.1 Khái quát về Kế toán Thuế hiện đại}
    \begin{itemize}
        \item Kế toán thuế là một quy trình đòi hỏi sự chính xác tuyệt đối và cập nhật liên tục.
        \item Hệ thống văn bản pháp luật thuế thường xuyên thay đổi và phức tạp.
        \item Các doanh nghiệp đối mặt với áp lực lớn về thời hạn nộp báo cáo và nghĩa vụ thuế.
    \end{itemize}
\end{frame}


\begin{frame}{1.2 Thách thức của Kế toán Thuế truyền thống}
    \begin{itemize}
        \item Khối lượng dữ liệu chứng từ (hóa đơn, biên lai) cần xử lý thủ công khổng lồ.
        \item Dễ xảy ra sai sót con người trong quá trình nhập liệu và đối chiếu.
        \item Tốn nhiều thời gian tra cứu và áp dụng các thông tư, nghị định mới.
    \end{itemize}
\end{frame}


\begin{frame}{1.3 Sự can thiệp của Trí tuệ Nhân tạo (AI)}
    \begin{itemize}
        \item AI mang đến giải pháp \textbf{tự động hóa} các quy trình thủ công.
        \item Cung cấp khả năng \textbf{phân tích dữ liệu lớn} (Big Data) để tìm ra các rủi ro tiềm ẩn.
        \item Chuyển dịch vai trò của kế toán viên từ "người nhập liệu" sang "nhà tư vấn chiến lược thuế".
    \end{itemize}
\end{frame}


\begin{frame}{1.4 Lợi ích cốt lõi khi áp dụng AI}
    \begin{itemize}
        \item \textbf{Tiết kiệm thời gian:} Giảm đến 80\% thời gian cho các công việc lặp đi lặp lại.
        \item \textbf{Tăng tính chính xác:} Loại bỏ rủi ro sai sót do yếu tố con người.
        \item \textbf{Đảm bảo tuân thủ:} AI tự động cập nhật và cảnh báo các thay đổi từ luật thuế.
    \end{itemize}
\end{frame}


\section{PHẦN 2: TỐI ƯU HÓA LẬP KẾ HOẠCH THUẾ VỚI AI}


\begin{frame}{2.1 Lập kế hoạch thuế là gì?}
    \begin{itemize}
        \item Là quá trình phân tích chiến lược tài chính nhằm tối thiểu hóa nghĩa vụ thuế.
        \item Vẫn phải đảm bảo tuân thủ nghiêm ngặt các quy định pháp luật.
        \item Mục tiêu là tối đa hóa lợi ích kinh tế cho doanh nghiệp (Tax Planning Optimization).
    \end{itemize}
\end{frame}


\begin{frame}{2.2 AI trong Tối ưu hóa Lập kế hoạch Thuế}
    \begin{itemize}
        \item AI phân tích dữ liệu tài chính, các quy định thuế và mục tiêu kinh doanh.
        \item Đề xuất các phương án lập kế hoạch thuế để tối đa hóa lợi ích.
        \item Xem xét đa chiều các yếu tố: khoản giảm trừ (deductions), tín dụng thuế (credits) và các trường hợp miễn giảm (exemptions).
    \end{itemize}
\end{frame}


\begin{frame}{2.3 Phân tích kịch bản Thuế (Scenario Analysis)}
    \begin{itemize}
        \item AI có khả năng tạo ra nhiều kịch bản (Tốt/Xấu/Cơ sở) cho quyết định kinh doanh.
        \item Dự báo dòng tiền thuế phải nộp trong tương lai dựa trên dữ liệu lịch sử.
        \item Giúp Ban giám đốc lựa chọn phương án có mức rủi ro và chi phí thuế thấp nhất.
    \end{itemize}
\end{frame}


\begin{frame}{2.4 Cập nhật Thông tin Thuế theo thời gian thực (Real-time Tax Insights)}
    \begin{itemize}
        \item AI liên tục theo dõi các thay đổi trong luật thuế, quy định và yêu cầu của từng ngành.
        \item Phân tích khối lượng lớn dữ liệu liên quan đến thuế để phát ra cảnh báo tức thì.
        \item Đảm bảo chiến lược thuế luôn được cập nhật và phù hợp với các quy định mới nhất.
    \end{itemize}
\end{frame}


\begin{frame}{2.5 Tính toán Thuế Tự động (Automated Tax Calculations)}
    \begin{itemize}
        \item Hệ thống AI phân tích dữ liệu tài chính, nhận diện các giao dịch chịu thuế.
        \item Tự động áp dụng các mức thuế suất và quy tắc tính thuế tương ứng.
        \item \textbf{Hình minh họa tính toán thuế:}
    \end{itemize}
    \begin{center}\includegraphics[width=0.7\textwidth,height=0.6\textheight,keepaspectratio]{images/Day\_11/docx\_img\_1.png}\end{center}
\end{frame}


\begin{frame}{2.6 Ưu điểm của Tính toán Tự động}
    \begin{itemize}
        \item Giảm thiểu rủi ro sai sót (human error).
        \item Đảm bảo tính chính xác cao trong quá trình tính thuế.
        \item Giải phóng thời gian quý báu của kế toán viên cho các nhiệm vụ mang tính chiến lược hơn.
    \end{itemize}
\end{frame}


\section{PHẦN 3: ỨNG DỤNG AI HỖ TRỢ KÊ KHAI \& TUÂN THỦ THUẾ}


\begin{frame}{3.1 Tuân thủ Thuế Tự động (Automated Tax Compliance)}
    \begin{itemize}
        \item AI hỗ trợ chuẩn bị và nộp các tờ khai thuế hoàn toàn tự động.
        \item Tích hợp trực tiếp với phần mềm kế toán để trích xuất dữ liệu tài chính liên quan.
        \item Tự động tạo các mẫu biểu thuế chính xác theo chuẩn mực của cơ quan thuế.
    \end{itemize}
\end{frame}


\begin{frame}{3.2 Lợi ích của Kê khai Tự động}
    \begin{itemize}
        \item Giảm rủi ro điền sai thông tin trên tờ khai.
        \item Đảm bảo tuân thủ đúng thời hạn, tránh các khoản phạt nộp chậm.
        \item Tối ưu hóa nguồn lực nhân sự của phòng kế toán.
    \end{itemize}
\end{frame}


\begin{frame}{3.3 Quản lý Tài liệu Thuế (Tax Document Management)}
    \begin{itemize}
        \item AI đơn giản hóa việc quản lý tài liệu thông qua số hóa và tổ chức hồ sơ.
        \item Sử dụng công nghệ Nhận dạng Ký tự Quang học (OCR) để đọc dữ liệu.
        \item Tự động trích xuất thông tin từ hóa đơn, biên lai và báo cáo tài chính.
    \end{itemize}
\end{frame}


\begin{frame}{3.4 Khả năng truy xuất và tìm kiếm tài liệu}
    \begin{itemize}
        \item Việc tìm kiếm một chứng từ cũ phục vụ thanh tra thuế diễn ra trong vài giây.
        \item Hệ thống AI gắn thẻ (tagging) tự động dựa trên nội dung chứng từ.
        \item Đảm bảo tính toàn vẹn và an toàn của kho lưu trữ số.
    \end{itemize}
\end{frame}


\begin{frame}{3.5 AI trong kê khai Thuế GTGT và TNCN}
    \begin{itemize}
        \item \textbf{Thuế GTGT:} Tự động đối chiếu hóa đơn đầu vào - đầu ra, loại bỏ hóa đơn không hợp lệ.
        \item \textbf{Thuế TNCN:} Tự động tập hợp thu nhập từ nhiều nguồn, tính toán mức giảm trừ gia cảnh.
        \item Gợi ý mức đóng thuế tối ưu cho người lao động và doanh nghiệp.
    \end{itemize}
\end{frame}


\begin{frame}{3.6 So sánh: Kế toán truyền thống vs. Kế toán AI}
    \begin{itemize}
        \item \textbf{Truyền thống:} Xử lý hàng ngàn hóa đơn thủ công cuối tháng, dễ sót hóa đơn.
        \item \textbf{AI:} Hóa đơn được xử lý ngay khi phát sinh, tự động kiểm tra tính hợp lệ qua kết nối API.
    \end{itemize}
\end{frame}


\section{PHẦN 4: PHÁT HIỆN GIAN LẬN \& QUẢN TRỊ RỦI RO THUẾ}


\begin{frame}{4.1 Đánh giá và Giảm thiểu Rủi ro (Risk Assessment)}
    \begin{itemize}
        \item AI hỗ trợ nhận diện các rủi ro thuế tiềm ẩn.
        \item Giảm thiểu khả năng doanh nghiệp rơi vào tình trạng không tuân thủ.
        \item Phân tích sâu các mẫu giao dịch và yêu cầu pháp lý để cắm cờ (flag) cảnh báo.
    \end{itemize}
\end{frame}


\begin{frame}{4.2 Tính chủ động trong Quản trị Rủi ro}
    \begin{itemize}
        \item Hệ thống không đợi đến lúc quyết toán mới báo lỗi.
        \item Đề xuất các chiến lược giảm thiểu rủi ro ngay khi giao dịch vừa được ghi nhận.
        \item Tránh được các khoản phạt nặng nề hoặc bị thanh tra đột xuất.
    \end{itemize}
\end{frame}


\begin{frame}{4.3 Phát hiện Gian lận Thuế (Tax Fraud Detection)}
    \begin{itemize}
        \item AI phân tích dữ liệu tài chính để nhận diện những bất thường (anomalies).
        \item So sánh các giao dịch hiện tại với dữ liệu lịch sử và mô hình hành vi.
        \item Cảnh báo các hành vi có dấu hiệu thao túng, chuyển giá hoặc trốn thuế.
    \end{itemize}
\end{frame}


\begin{frame}{4.4 Tầm quan trọng của Phát hiện Gian lận sớm}
    \begin{itemize}
        \item Ngăn chặn thất thoát tài sản của doanh nghiệp.
        \item Bảo vệ danh tiếng và uy tín của công ty trên thị trường.
        \item Đảm bảo tuân thủ tuyệt đối các quy định của cơ quan thuế.
    \end{itemize}
\end{frame}


\begin{frame}{4.5 Hỗ trợ Thanh tra Thuế (Audit Support)}
    \begin{itemize}
        \item AI tổ chức và truy xuất dữ liệu tài chính một cách hiệu quả nhất.
        \item Hỗ trợ cung cấp hồ sơ kiểm toán bằng các thuật toán Machine Learning.
        \item Xác định sự thiếu nhất quán trong sổ sách trước khi cơ quan thuế phát hiện.
        \item \textbf{Hỗ trợ thanh tra:}
    \end{itemize}
    \begin{center}\includegraphics[width=0.7\textwidth,height=0.6\textheight,keepaspectratio]{images/Day\_11/docx\_img\_2.jpeg}\end{center}
\end{frame}


\begin{frame}{4.6 Audit Trail (Dấu vết Kiểm toán)}
    \begin{itemize}
        \item Mọi chỉnh sửa, thao tác trên chứng từ đều được AI lưu lại.
        \item Dễ dàng truy vết ai là người thay đổi số liệu và thay đổi vì lý do gì.
        \item Nâng cao tính minh bạch và sức thuyết phục khi giải trình với cơ quan chức năng.
    \end{itemize}
\end{frame}


\section{PHẦN 5: ỨNG DỤNG AI VỚI LUẬT QUẢN LÝ THUẾ \& THÔNG TƯ 200}


\begin{frame}{5.1 Xử lý văn bản pháp luật bằng AI}
    \begin{itemize}
        \item Các thông tư (như Thông tư 200) và văn bản luật rất dài và phức tạp.
        \item AI (như ChatGPT) có thể "đọc" hàng trăm trang tài liệu trong vài giây.
        \item Khả năng xử lý ngôn ngữ tự nhiên giúp AI hiểu đúng ngữ cảnh pháp lý.
    \end{itemize}
\end{frame}


\begin{frame}{5.2 Kỹ thuật Tóm tắt Văn bản Phức tạp}
    \begin{itemize}
        \item Kế toán viên có thể sử dụng các "Prompt" (Câu lệnh) để yêu cầu AI tóm tắt.
        \item Ví dụ: "Tóm tắt các điều kiện ghi nhận doanh thu theo Thông tư 200".
        \item Kết quả trả về là các gạch đầu dòng ngắn gọn, dễ hiểu và đi thẳng vào trọng tâm.
    \end{itemize}
\end{frame}


\begin{frame}{5.3 Nhận diện rủi ro từ Thông tư 200}
    \begin{itemize}
        \item AI có thể so sánh dữ liệu hạch toán thực tế với quy định trong Thông tư 200.
        \item Phát hiện các bút toán không đúng bản chất theo hướng dẫn của Chế độ kế toán.
        \item Đảm bảo tính tuân thủ ngay từ khâu ghi sổ (bookkeeping).
    \end{itemize}
\end{frame}


\begin{frame}{5.4 Phân tích Điều kiện Khấu trừ Thuế}
    \begin{itemize}
        \item Đưa một trường hợp chi phí cụ thể vào AI và hỏi: "Chi phí này có được khấu trừ thuế TNDN không?"
        \item AI sẽ quét các quy định trong Luật Quản lý Thuế và phản hồi kèm theo trích dẫn.
        \item Giúp kế toán viên tự tin hơn trong các quyết định hạch toán.
    \end{itemize}
\end{frame}


\begin{frame}{5.5 Cập nhật chênh lệch giữa Kế toán và Thuế}
    \begin{itemize}
        \item Kế toán tài chính tuân theo Thông tư 200, trong khi Kế toán thuế tuân theo Luật Thuế.
        \item AI hỗ trợ tự động bóc tách và theo dõi các khoản chênh lệch tạm thời/vĩnh viễn.
        \item Hỗ trợ lập Bảng kê khai quyết toán thuế TNDN nhanh chóng.
    \end{itemize}
\end{frame}


\begin{frame}{5.6 Tính năng "Trợ lý ảo Pháp lý"}
    \begin{itemize}
        \item Có thể huấn luyện một Chatbot AI nội bộ chứa dữ liệu là toàn bộ hệ thống Luật Thuế Việt Nam.
        \item Kế toán viên có thể đặt câu hỏi bằng tiếng Việt thông thường.
        \item Tiết kiệm chi phí thuê chuyên gia tư vấn luật cho các vấn đề cơ bản.
    \end{itemize}
\end{frame}


\section{PHẦN 6: GIỚI HẠN PHÁP LÝ \& ĐẠO ĐỨC NGHỀ NGHIỆP}


\begin{frame}{6.1 Rủi ro "Ảo giác" (Hallucinations) của AI}
    \begin{itemize}
        \item AI tạo sinh (như ChatGPT) có thể tự bịa ra các điều luật không có thật.
        \item Nguy hiểm nếu kế toán viên áp dụng mù quáng mà không kiểm chứng.
        \item Luôn phải có bước xác minh (verify) bằng các nguồn tài liệu chính thống.
    \end{itemize}
\end{frame}


\begin{frame}{6.2 Bảo mật Dữ liệu Tài chính}
    \begin{itemize}
        \item Việc đưa dữ liệu công ty (doanh thu, lợi nhuận, lương) lên các nền tảng AI công cộng có thể gây rò rỉ.
        \item Các doanh nghiệp cần chính sách bảo mật dữ liệu nghiêm ngặt.
        \item Ưu tiên sử dụng các hệ thống AI nội bộ (On-premise) hoặc có cam kết bảo mật cấp doanh nghiệp (Enterprise Tier).
    \end{itemize}
\end{frame}


\begin{frame}{6.3 Trách nhiệm Giải trình (Accountability)}
    \begin{itemize}
        \item Khi AI tính sai thuế, cơ quan thuế sẽ phạt doanh nghiệp, không phải phạt AI.
        \item Người kế toán viên vẫn là người chịu trách nhiệm cuối cùng (Human-in-the-loop).
        \item AI chỉ là công cụ hỗ trợ (Copilot), không phải người thay thế hoàn toàn (Autopilot).
    \end{itemize}
\end{frame}


\begin{frame}{6.4 Đạo đức Nghề nghiệp Kế toán}
    \begin{itemize}
        \item Tuân thủ nguyên tắc khách quan và trung thực.
        \item Không sử dụng AI để tìm kẽ hở trốn thuế một cách phi pháp.
        \item Đảm bảo ứng dụng AI nhằm tăng cường tính minh bạch, không phải để che giấu sai phạm.
    \end{itemize}
\end{frame}


\section{PHẦN 7: TỔNG KẾT \& THẢO LUẬN}


\begin{frame}{7.1 Tóm tắt Bài học}
    \begin{itemize}
        \item AI mang lại bước tiến lớn trong tối ưu hóa kế hoạch thuế và tuân thủ.
        \item Tự động hóa tính toán, kê khai và chuẩn bị hồ sơ thanh tra.
        \item Ứng dụng AI giúp tra cứu, tóm tắt Thông tư 200 và Luật Thuế dễ dàng.
        \item Luôn đề cao bảo mật dữ liệu và kiểm chứng thông tin từ AI.
    \end{itemize}
\end{frame}


\begin{frame}{7.2 Định hướng thực hành (Buổi 11 - Thực hành)}
    \begin{itemize}
        \item Học viên sẽ được thực hành:
        \item Đưa một đoạn Thông tư Thuế phức tạp vào ChatGPT để yêu cầu tóm tắt.
        \item Sử dụng Prompt liệt kê các điều kiện khấu trừ thuế từ tình huống thực tế.
        \item Kiểm tra tính đúng đắn và xử lý rủi ro "ảo giác" của AI.
    \end{itemize}
\end{frame}


\begin{frame}{7.3 Thảo luận mở}
    \begin{itemize}
        \item \textbf{Câu hỏi:} "Bạn nghĩ rào cản lớn nhất khi doanh nghiệp Việt Nam áp dụng AI vào Kế toán Thuế là gì?"
        \item \textbf{Q\&A:} Giải đáp các thắc mắc của học viên.
    \end{itemize}
\end{frame}


\begin{frame}{7.4 Tài liệu Đọc thêm}
    \begin{itemize}
        \item *Chapter 9: AI-Based Tax Planning and Compliance (Charnelle Gibson)*
        \item *Luật Quản lý thuế, Thông tư 200 (Bản Tiếng Việt)*
        \item *Sổ tay Prompt Kế toán (Scott Dell)*
    \end{itemize}
\end{frame}


\begin{frame}{CẢM ƠN CÁC BẠN ĐÃ LẮNG NGHE!}
    \begin{itemize}
        \item Chúc các bạn ứng dụng thành công AI vào Kế toán Thuế!
    \end{itemize}
\end{frame}


\end{document}